\chapter{The Compiler to C}\label{ch:dylanc}
\sloppy

The second implementation in this system is a compiler. It reads a whole
Dylan program, translates it into C, and runs a C compiler to make a
native executable. The executable needs nothing from SBCL at run time: it
is an ordinary program linked with a small runtime library written in C,
with a garbage collector --- the Memory Pool System, or by choice the
Boehm-Demers-Weiser collector (Section~\ref{sec:dylanc-gc}) --- and
with the system's thread library.

This chapter is about using the compiler: the \cmd{bin/dylanc} command
and every one of its options, the files it writes, the cache it keeps,
the kinds of build it can make, how to call it from Common Lisp, what a
compiled program does when something goes wrong, and what a program must
not do because it is compiled whole and in advance. How the compiler
works inside is the subject of Chapter~\ref{ch:internals}, and how to
make compiled programs fast is the subject of
Chapter~\ref{ch:performance}.

\section{A first build}\label{sec:dylanc-first}

\index{dylanc@\cmd{dylanc}}
A program is one or more files of Dylan source, conventionally with the
extension \file{.dyl}. Here is the smallest useful one:

\begin{dylan}
(define-method greet ((name <string>))
  (format #t "Hello, ~a!~%" name))
(greet "world")
\end{dylan}

Compile it and run the result:

\begin{shell}
$ bin/dylanc hello.dyl
$ ./hello
Hello, world!
\end{shell}

\cmd{bin/dylanc} is a short shell script that runs
\file{compiler/dylanc.lisp} under \cmd{sbcl -{}-script}, so SBCL must be
on your \code{PATH}; the script can be run from any directory, and can be
linked into a directory on your \code{PATH}. The executable is named
after the first file, less its \file{.dyl}, and is written beside that
file: \cmd{bin/dylanc src/hello.dyl} makes \file{src/hello}.

The top-level forms of the program run in order, from the first file to
the last, when the executable starts. There is no \code{main} function to
write: the program \emph{is} its top-level forms, as it would be if you
typed them into the interactive implementation's read-eval-print loop. A
top-level form's value is not printed; use \code{print}, \code{format}
and the other printing functions (Chapter~\ref{ch:printing}) to produce
output. When the last top-level form has run, the program exits with
status~0.

\paragraph{Several files.}
Give several files to make one program from them all:

\begin{shell}
$ bin/dylanc -o app utilities.dyl shapes.dyl main.dyl
\end{shell}

The files are read in the order given, and their top-level forms run in
that order. A file whose first line is a \code{module:} header (design
note~33; see Chapter~\ref{ch:modules}) is compiled in the module it
names; any other file is compiled in \code{dylan-user}. Every
\code{define-module} form in every file is known before any other form is
compiled, so a module may be defined after the files that use it.

\paragraph{What a first build costs.}
The compiler is SBCL code, compiled into \file{compiler/fasl/} the first
time it is used, and the runtime library is compiled by \cmd{make} the
first time a build needs it. After that, a small program builds in about a
second and a half on a current machine. The first build with a new
combination of C compiler and options takes several seconds longer,
because the library's C must be compiled for it
(Section~\ref{sec:dylanc-cache}).

\section{The options}\label{sec:dylanc-options}

\index{dylanc@\cmd{dylanc}!options}
\cmd{bin/dylanc -h} lists the options:

\begin{shell}
Usage: dylanc [options] file.dyl ...

Compile the Dylan files, in order, into one executable.

Options:
  -o FILE          the executable (default: the first file's name, less .dyl)
  -O0 -O1 -O2 -O3  how much the C compiler optimizes (default -O2)
  -Os              a small executable: optimized for size, unused code dropped
                   (the library compiled with the program, as --whole)
  --whole          compile the library with the program, not apart (the
                   library's C, compiled apart, is compiled once and kept)
  --lto            optimize the program and the runtime library together when
                   linking (link-time optimization)
  -g, --debug      for debugging: no optimization, debugging information, and
                   every inline operation checking its arguments when run
  -k, --keep-c     keep the C: FILE.c, FILE-MODULE.c for each module, FILE.h
  -C, --c-only     write the C only (FILE.c and the rest); do not compile it
  -E, --expand     write each use of the program's macros and its expansion;
                   no C
  --trampoline     tail calls by trampolines, for C compilers without musttail
  --cc CC          the C compiler (default: $CC, or cc)
  --no-lines       no #line directives pointing to the Dylan source
  --sanitize KIND  for testing: a debug build with the sanitizers, KIND
                   address (AddressSanitizer and UndefinedBehaviorSanitizer)
                   or thread (ThreadSanitizer)
  --gc NAME        the garbage collector: mps (the default), the Memory Pool
                   System, generational and moving; or boehm, the installed
                   Boehm-Demers-Weiser collector (libgc)
  --gc-stress      for testing: a collection every DYLAN_GC_STRESS (default
                   100) allocations
  -h, --help       this
\end{shell}

Options and file names may be mixed in any order. Each option is
described below, with the occasions for using it.

\begin{description}
\item[\cmd{-o} \meta{file}] \index{dylanc@\cmd{dylanc}!-o@\cmd{-o}}
  Names the executable. Without it the executable is named after the
  first file. The C files, when they are kept, are named after the
  executable: \cmd{-o build/app} writes \file{build/app.c} and the rest.

\item[\cmd{-O0}, \cmd{-O1}, \cmd{-O2}, \cmd{-O3}]
  How hard the C compiler optimizes. The default, \cmd{-O2}, is right for
  almost every use. \cmd{-O0} compiles fastest, which matters only for
  large programs while you are developing them; measure before assuming
  that \cmd{-O3} helps, since it makes larger code. An
  ordinary build is also compiled with \cmd{-g}, so a debugger can show
  the Dylan source even in an optimized program (see
  Section~\ref{sec:dylanc-debug}). The Dylan compiler's own optimizations
  (Chapter~\ref{ch:internals}) are always on; these options affect only
  the C compiler.

\item[\cmd{-Os}] \index{size builds}
  Makes a small executable. The library is compiled with the program,
  and only the parts of it the program can reach are kept (\term{tree
  shaking}); the C is compiled with \cmd{-Os}, each function and datum in
  its own section, the linker drops the sections nothing uses, and the
  executable is stripped of its symbols. Use it for programs you ship or
  copy elsewhere, when size matters more than build time. Section~\ref{sec:dylanc-size}
  shows what it gains.

\item[\cmd{-{}-whole}]
  Compiles the library with the program, as \cmd{-Os} does, and shakes it
  the same way, but optimizes for speed as usual. Normally the library is
  compiled apart from the program, once, and kept
  (Section~\ref{sec:dylanc-cache}); with \cmd{-{}-whole} its C is compiled
  afresh with every build. The result is smaller than an ordinary build,
  and a build takes about twice as long. It is mostly a check that the
  separately compiled library and the shaken one agree, which the test
  suites make (Appendix~\ref{app:testing}).

\item[\cmd{-{}-lto}] \index{link-time optimization}
  Link-time optimization: the C compiler optimizes the program and the C
  runtime library together when it links them, so that small runtime
  functions can be inlined into the program. The runtime library is built
  for this the first time it is asked for, with \cmd{-O2 -g -flto}, in
  \file{runtime/build/lto-\meta{cc}/}. The effect on speed is usually
  small; try it on your own program and measure. It is ignored by
  \cmd{-g} and by the testing builds below.

\item[\cmd{-g}, \cmd{-{}-debug}]
  A build for debugging: the C is compiled with \cmd{-O0 -g
  -DDY\_DEBUG}, and the compiler trusts none of the types it infers, so
  every operation it would have done inline tests its arguments when the
  program runs. Section~\ref{sec:dylanc-debug} has more.

\item[\cmd{-k}, \cmd{-{}-keep-c}]
  Keeps the C files, which are otherwise deleted once the executable is
  made. Use it to read what the compiler wrote
  (Section~\ref{sec:internals-reading}), or to see what a debugger is
  showing you.

\item[\cmd{-C}, \cmd{-{}-c-only}]
  Writes the C files and stops: no C compiler is run and no executable is
  made. Use it to compile the C elsewhere, with a build system of your
  own, or on a machine without SBCL. Section~\ref{sec:dylanc-files}
  shows the command that compiles the result by hand.

\item[\cmd{-E}, \cmd{-{}-expand}]
  Writes each use of the program's own macros, with the place it occurs
  and what it expands into, and stops; no C is written. Use it to see
  what your macros do (Chapter~\ref{ch:extensions}).

\item[\cmd{-{}-trampoline}] \index{trampoline}
  Makes tail calls with trampolines instead of the \code{musttail}
  attribute. Every Dylan call in tail position must run in constant
  stack space, and C does not promise that; GCC~15 and later and Clang~13
  and later do, through \code{musttail}. With an older or different C
  compiler, \cmd{-{}-trampoline} makes each tail call return to a loop
  that makes the call, which any C11 compiler can do. Programs built this
  way are correct but slower: a call-heavy benchmark (\code{tak}) took
  about 1.7 times as long. Without this option, a C compiler that lacks
  \code{musttail} stops with the error ``This C compiler cannot guarantee
  tail calls (musttail): define DY\_TRAMPOLINE.''

\item[\cmd{-{}-cc} \meta{cc}]
  The C compiler to run, such as \cmd{gcc} or \cmd{clang}, or a full path.
  Without it, the compiler is the value of the \code{CC} environment
  variable, or \cmd{cc} if that is not set. It is worth trying both GCC
  and Clang on a program whose speed matters.

\item[\cmd{-{}-no-lines}]
  Leaves out the \code{\#line} directives that point the C back to the
  Dylan source (Section~\ref{sec:dylanc-files}). C compilers' messages,
  and debuggers, then refer to the C files themselves. Use it when you
  want to step through the generated C rather than the Dylan.

\item[\cmd{-{}-gc} \meta{name}] \index{garbage collection!choosing}
  The garbage collector the program is built with: \cmd{mps}, the Memory
  Pool System, the default, which collects by generations and moves
  objects; or \cmd{boehm}, the Boehm-Demers-Weiser collector, which must
  be installed. See Section~\ref{sec:dylanc-gc}. The environment variable
  \code{DYLAN\_GC=boehm} makes \cmd{boehm} the default.

\item[\cmd{-{}-sanitize} \meta{kind}]
  For testing the compiler and runtime: a debug build with
  \meta{kind} \cmd{address} (AddressSanitizer and
  UndefinedBehaviorSanitizer) or \cmd{thread} (ThreadSanitizer). See
  Section~\ref{sec:dylanc-testing-builds}.

\item[\cmd{-{}-gc-stress}]
  For testing: a garbage collection every few allocations. See
  Section~\ref{sec:dylanc-testing-builds}.

\item[\cmd{-h}, \cmd{-{}-help}]
  Prints the summary above and exits with status 0.
\end{description}

When options conflict, the more specific build wins: a testing build
(\cmd{-{}-sanitize}, \cmd{-{}-gc-stress}) uses its own C flags whatever
else is given; otherwise \cmd{-g} wins over \cmd{-Os}, and \cmd{-Os} over
\cmd{-O}\meta{n}.

\section{The files the compiler writes}\label{sec:dylanc-files}

\index{C files}
For a program whose executable is \meta{name}, the compiler writes these
C files:

\begin{description}
\item[\file{\meta{name}.c}] The program's static data --- every class,
  symbol, keyword and constant, laid out as the runtime expects --- with
  the object system's functions (the dispatch function of each generic
  function, slot accessors, and the makers \code{make} calls), and
  \code{main}.
\item[\file{\meta{name}-\meta{module}.c}] The functions written in each
  module, one file for each module. A program written in
  \code{dylan-user} has \file{\meta{name}-dylan-user.c}; a program with
  modules \code{counting} and \code{app} has
  \file{\meta{name}-counting.c} and \file{\meta{name}-app.c} too.
\item[\file{\meta{name}-dylan.c}] The library's functions: the module
  \code{dylan}, written in Dylan in \file{runtime/dylan/} and compiled
  with every program. This is the largest file by far, about a megabyte.
\item[\file{\meta{name}.h}] What the files share: a declaration of every
  object \file{\meta{name}.c} defines and a prototype of every function.
\end{description}

The C is linked with \file{runtime/build/libdylan.a}, the runtime library
compiled from \file{runtime/c/} with the collector in it, and with
\cmd{-lm -lpthread} (with \cmd{-{}-gc boehm}, the library is
\file{runtime/build/boehm-\meta{cc}/libdylan.a} and \cmd{-lgc} is
linked too). With \cmd{-k} the files stay after the build:

\begin{shell}
$ bin/dylanc -k hello.dyl
$ ls hello*
hello  hello.c  hello.dyl  hello-dylan.c  hello-dylan-user.c  hello.h
\end{shell}

\paragraph{\code{\#line} directives.} \index{\#line directives}
The generated C carries \code{\#line} directives that name the Dylan file
and line each statement came from, and after each function a directive
back to the C file itself. A C compiler's warnings, a debugger's
breakpoints and stack traces, and a profiler's reports therefore name
places in your Dylan source. Here is the start of the function
\file{hello-dylan-user.c} holds for \code{greet}:

\begin{ccode}
D fn1190_greet(D self, D a0, D a1, D a2, D a3, intptr_t argc) {
    (void)self; (void)a0; (void)a1; (void)a2; (void)a3; (void)argc;
    D v105609, v105610, v105611;
#line 1 "hello.dyl"
    if (DY_UNLIKELY(argc != 1)) dy_wrong_argument_count(self, argc);
    v105609 = a0;
#line 2 "hello.dyl"
    ...
\end{ccode}

\paragraph{Compiling the C yourself.}
With \cmd{-C}, the compiler writes the C and stops. The files can then be
compiled with the C compiler's options and the runtime library's
directory:

\begin{shell}
$ bin/dylanc -C hello.dyl
$ cc -O2 -std=c11 -pthread -DDY_MPS \
     -I$DYLAN/runtime/c -I$DYLAN/runtime/mps/code \
     hello.c hello-dylan.c hello-dylan-user.c \
     $DYLAN/runtime/build/libdylan.a -lm -lpthread -o hello
$ ./hello
Hello, world!
\end{shell}

Here \code{\$DYLAN} stands for the root of the source tree. The runtime
library \file{libdylan.a} must exist; \cmd{make -C runtime} makes it, and
any build with \cmd{dylanc} makes it too. Add \cmd{-DDY\_TRAMPOLINE} to
the C compiler's options, and link \file{libdylan-trampoline.a} instead,
if your C compiler lacks \code{musttail}. For the Boehm collector, leave
out \cmd{-DDY\_MPS} and the MPS include directory, link
\file{runtime/build/boehm-cc/libdylan.a} (made by any
\cmd{dylanc -{}-gc boehm} build), and add \cmd{-lgc}. The C is C11 with POSIX
threads; \file{runtime/c/dylan.h} checks what it assumes (64-bit
pointers, arithmetic right shifts) with \code{\_Static\_assert}, and
checks that the runtime it is compiled against is the version the
compiler wrote for.

\section{The library's cache}\label{sec:dylanc-cache}

\index{library cache}
Every program carries the whole Dylan library, about a megabyte of C. To
avoid compiling that with every build, the compiler writes the library's
C so that it does not depend on the program: \file{\meta{name}-dylan.c}
is the same text for almost every program. The driver compiles it once
for each C compiler and set of options and keeps the object in
\file{runtime/build/library/}, named for a digest (MD5) of the C text,
the C compiler's name and \cmd{-{}-version} output, the options, and the
runtime headers \file{dylan.h} and \file{primitive.h}. A later build that
would compile the same text the same way uses the object instead. This
halves the time of a small build.

The newest 40 objects are kept, with the C each was compiled from; older
ones are deleted. The cache is safe to delete at any time
(\cmd{rm -r runtime/build/library}); it is rebuilt as needed. Builds
running at the same time share it safely.

A program can change the library's C --- for example by defining a method
that replaces one of the library's own. Then the digest differs and the
library is simply compiled for that build. \cmd{-Os} and
\cmd{-{}-whole} do not use the cache at all: they compile the library
with the program, keeping only what the program can reach.

\section{Size}\label{sec:dylanc-size}

These are the sizes of the \code{greet} program above, built with GCC~16
on x86-64 Linux, as reported by \cmd{size} (the code, the initialized
data and the uninitialized data, excluding the shared C and collector
libraries):

\begin{center}
\begin{tabular}{lrr}
\toprule
Build & text + data (KB) & file size (KB) \\
\midrule
default (\cmd{-O2}) & 463 & 2765 \\
\cmd{-{}-whole} & 307 & 1798 \\
\cmd{-{}-lto} & 334 & 2374 \\
\cmd{-Os} & 221 & 231 \\
\bottomrule
\end{tabular}
\end{center}

The file sizes of all but the \cmd{-Os} build include debugging
information. Hello world built \cmd{-Os} is about 220~KB of code and data;
most of it is library code that a call of a generic function such as
\code{print-object}, \code{make} or \code{as} could reach.
Chapter~\ref{ch:performance} says more about size.

\section{Debug builds}\label{sec:dylanc-debug}

\index{debugging}
A program built with \cmd{-g} differs from an ordinary one in three ways:

\begin{itemize}
\item The C is compiled with \cmd{-O0 -g}, so a debugger can show every
  variable and step through every line.
\item The compiler trusts none of the classes it has inferred. Where an
  ordinary build adds two integers with no test because it has proved
  them integers, a debug build tests them anyway. A debug build therefore
  catches a mistake in the compiler's reasoning as an error, rather than
  as a crash or a wrong answer.
\item \code{DY\_DEBUG} is defined, which makes the runtime's accessors
  check the representation of each value they take apart: that what is
  read as a pair \emph{is} a pair, and so on. A violation stops the
  program with a message naming what was expected.
\end{itemize}

None of these change what a correct program does: a debug build reports
the same Dylan errors, in the same way, as an ordinary one. Debug builds
are for finding bugs in the compiler or runtime, and for using a
debugger. Because of the \code{\#line} directives, a debugger works in
Dylan terms --- breakpoints can be set at lines of the Dylan file. Given
this \file{sum.dyl},

\begin{dylan}
(define-method sum-to ((n <integer>))
  (for ((i 0 (+ i 1))
        (s 0 (+ s i)))
       ((> i n) s)))
(print (sum-to 100))
\end{dylan}

\noindent a debug build stops where it is asked to:

\begin{shell}
$ bin/dylanc -g -o sum sum.dyl
$ gdb -batch -ex "break sum.dyl:3" -ex run sum
...
Thread 1 "sum" hit Breakpoint 1, fn1190_sum_to (self=4692864, a0=201,
    a1=0, a2=0, a3=0, argc=1) at /home/me/sum.dyl:3
3               (s 0 (+ s i)))
\end{shell}

The names of the C functions are those the compiler gave the Dylan
functions (\code{fn1190\_sum\_to} is the method \code{sum-to}), and the
values are Dylan values in their machine representation: \code{a0=201}
is the integer 100, stored shifted left one bit with the low bit set.
Chapter~\ref{ch:internals} explains the representation.

\section{Testing builds}\label{sec:dylanc-testing-builds}

Two options build programs for testing the compiler and runtime rather
than for use. Each builds its own variant of the runtime library, the
first time it is needed, in \file{runtime/build/\meta{variant}-\meta{cc}/},
with the same C flags as the program.

\begin{description}
\item[\cmd{-{}-sanitize address}] Compiles with \cmd{-O1 -g
  -fno-omit-frame-pointer -fsanitize=address,undefined -DDY\_DEBUG}, as a
  debug build: AddressSanitizer reports any read or write outside an
  object, and UndefinedBehaviorSanitizer any undefined behavior in C,
  such as a signed overflow or a misaligned access.
\item[\cmd{-{}-sanitize thread}] Compiles with \cmd{-O1 -g
  -fsanitize=thread -DDY\_DEBUG}: ThreadSanitizer reports data races
  between threads.
\item[\cmd{-{}-gc-stress}] Compiles with \cmd{-O2 -g -DDY\_GC\_STRESS}:
  the program forces a garbage collection every \meta{n} allocations,
  where \meta{n} is the value of the environment variable
  \code{DYLAN\_GC\_STRESS}, or 100. A pointer the collector cannot see
  --- a bug in the compiler or runtime --- shows up quickly as a crash or
  a wrong answer. Such a program runs very slowly.
\end{description}

\begin{shell}
$ bin/dylanc --gc-stress -o hgc hello.dyl
$ DYLAN_GC_STRESS=10 ./hgc
Hello, world!
\end{shell}

One of AddressSanitizer's checks is turned off by the runtime itself: its
detection of a frame used after its function returned, which keeps the
frames of functions apart from the stack. The collector finds the
references a program holds by scanning its stack, and would not see those
frames; an object referred to only from one would be freed, or moved,
while in use. The runtime supplies the option
\code{detect\_stack\_use\_after\_return=0} as the sanitizer's default,
so if you set \code{ASAN\_OPTIONS} yourself, leave that one off.

The test suites use these builds (Appendix~\ref{app:testing}). An
ordinary Dylan program has no need of them, but a program that misbehaves
in a way you cannot explain is worth running under them: what they find
is a bug to report.

\section{The garbage collector}\label{sec:dylanc-gc}

\index{garbage collection!choosing}\index{Memory Pool System}
\index{Boehm collector}
A compiled program's memory is managed by one of two collectors, chosen
when it is built.

\begin{description}
\item[\cmd{-{}-gc mps}] (the default) The Memory Pool System (MPS),
  Ravenbrook Limited's collector, written in the first place for
  Harlequin's Dylan. Its source is in the tree (\file{runtime/mps/},
  under its own BSD license) and is compiled into the runtime library;
  nothing need be installed. It is \emph{generational} and
  \emph{moving}: new objects go in a nursery, a collection of the nursery
  copies the few that survive out of it and leaves the rest behind
  without visiting them, and the older generations are collected rarely.
  The stacks and the program's static data are scanned conservatively,
  so an object they point to stays in place; the heap itself is scanned
  exactly. Each thread allocates by bumping a pointer in its own region,
  without locking.
\item[\cmd{-{}-gc boehm}] The Boehm-Demers-Weiser collector,
  \file{libgc}, which must be installed on the system
  (Section~\ref{sec:start-requirements}). It is \emph{conservative}: it
  treats any word that looks like a pointer as one, so it needs no
  knowledge of the program's objects, and objects never move. It
  collects by marking and sweeping the whole heap, in parallel on two
  threads, and allocates from per-thread free lists. It runs on more
  platforms than the MPS, and is the choice where the MPS has no port.
\end{description}

The MPS is the faster. A program that makes many short-lived objects
--- lists built and discarded, boxed floats, strings, conditions ---
gains most, two to three times on such programs, since a collection
costs in proportion to what survives, not to what was allocated; a
program that keeps a large heap alive gains from the old generations
being collected rarely; and a program that allocates little sees no
difference. The price is memory: a program that allocates heavily uses
about twice as much under the MPS, whose nursery is 128~MB and whose
heap may double between collections of everything, and
Section~\ref{sec:perf-gc} tells how to trade that back. The two
collectors are the same to the Dylan program, with one visible
exception: under the MPS an object's identity hash is a number given it
the first time it is hashed, not its address, since the address changes.

Every object of a program must be built for the same collector, and the
runtime library for each is kept apart: the MPS one is
\file{runtime/build/libdylan.a}, and the Boehm one is in
\file{runtime/build/boehm-\meta{cc}/}, with the variant and
\cmd{-{}-trampoline} in the name as for testing builds
(Section~\ref{sec:dylanc-testing-builds}). The testing builds work with
either collector. Under the MPS, \cmd{-{}-gc-stress} forces a
collection of every generation, and a ThreadSanitizer build uses the MPS
without its memory barrier, whose signal handling and the sanitizer's
would deadlock, so its collections there are not incremental. Setting
the environment variable \code{DYLAN\_GC=boehm} makes the Boehm
collector the default, which is how the test suites
(Appendix~\ref{app:testing}) are run with it. The MPS's behavior can be
tuned by environment variables, in Section~\ref{sec:perf-gc}. It has
been built here on Linux on x86-64 and macOS on Apple silicon, and has
ports to FreeBSD, to Windows with Visual C, and to Linux on ARM64.

\section{Calling the compiler from Lisp}\label{sec:dylanc-lisp}

\index{dylanc@\code{dylanc} package}
The compiler is a Common Lisp program, and it can be called from a Lisp
session. Load \file{compiler/load.lisp}, which loads the interactive
implementation's reader and the compiler, compiling any source newer
than its fasl; the package is \code{DYLAN-COMPILER}, nicknamed
\code{DYLANC}. Two functions do the work:

\begin{description}
\item[\code{(dylanc:build \meta{files} \meta{executable} \&key ...)}]
  Compiles the list of Dylan \meta{files} to C and on to the
  \meta{executable}. The keyword arguments correspond to the options of
  \cmd{dylanc}: \code{:optimize} (a string such as \code{"-O3"},
  default \code{"-O2"}), \code{:size}~\code{t} (\cmd{-Os}),
  \code{:whole}~\code{t}, \code{:lto}~\code{t}, \code{:debug}~\code{t},
  \code{:trampoline}~\code{t}, \code{:cc} (default \code{"cc"}; the
  \code{CC} environment variable is not consulted here),
  \code{:variant} (\code{:address}, \code{:thread} or
  \code{:gc-stress}), \code{:gc} (\code{:mps} or \code{:boehm}), and
  \code{:c-file}, the name of the main C file (default
  \meta{executable}\code{.c}). It answers \code{T} if the
  executable was made and \code{NIL} if not. Diagnostics are written to
  \code{*error-output*}. Unlike \cmd{dylanc}, it leaves the C files in
  place.
\item[\code{(dylanc:compile-files \meta{files} \meta{c-file})}]
  Writes the C only, as \cmd{-C} does. It answers the list of C files
  written, or \code{NIL} after reporting errors.
\end{description}

\begin{lisp}
(load "compiler/load.lisp")
(dylanc:build '("hello.dyl") "hello")              ; => T
(dylanc:build '("hello.dyl") "hello-small" :size t) ; => T
(dylanc:compile-files '("hello.dyl") "cf.c")
;; => (#P"cf.c" #P"cf-dylan.c" #P"cf-dylan-user.c")
\end{lisp}

Relative file names are taken relative to the Lisp's current directory
(\code{*default-pathname-defaults*}).

\section{Errors when compiling}\label{sec:dylanc-diagnostics}

\index{diagnostics}
The compiler reports problems as \term{diagnostics}, one to a line, on
the standard error stream, in the form \meta{file}\code{:}\meta{line}%
\code{:}\meta{column}\code{: }\meta{kind}\code{: }\meta{message}, where
\meta{kind} is \code{error} or \code{warning}:

\begin{shell}
$ cat r1.dyl
(print (undefined-thing 3))
$ bin/dylanc r1.dyl
r1.dyl:1:8: warning: undefined-thing is not defined (module dylan-user)
\end{shell}

A \textbf{warning} does not stop the build. The commonest is a variable
that is referred to but defined nowhere in the program; the program is
built, and fails only if the reference is reached (see
Section~\ref{sec:dylanc-runtime-errors}). Treat such a warning as an
error unless you know why it is there: it is usually a misspelling.

An \textbf{error} stops the build: no executable is made and
\cmd{dylanc} exits with status~1. Errors found while reading and
converting the source --- an unbalanced parenthesis, a malformed special
form, an integer too large --- are all reported in one run; the compiler
abandons the top-level form that has the error and goes on to the next.
Errors found later, when the program is analysed as a whole, stop at the
first. For example:

\begin{shell}
$ cat err2.dyl
(define f (method (x) (+ x 1))
(print (f 1))
$ bin/dylanc err2.dyl
err2.dyl:2:14: error: End of file inside a Dylan list.
\end{shell}

The position of a reading error is where the reader stopped, which for a
missing closing parenthesis is the end of the file; the unclosed list
began earlier.

\paragraph{Exit status.} \index{exit status!of dylanc@of \cmd{dylanc}}
\cmd{dylanc} exits with status 0 when it has made what was asked; 1 when
the program had errors or the C compiler failed (``\code{dylanc: the C
compiler failed.}''); and 2 for a bad command line --- an unknown option,
an option missing its value, no files, or a file that does not exist ---
after printing the usage summary.

\section{Errors when running}\label{sec:dylanc-runtime-errors}

\index{unhandled conditions}
A compiled program signals conditions as Chapter~\ref{ch:conditions}
describes, and handlers established with \code{handler-case} or
\code{handler-bind} see them as they would in the interactive
implementation. The runtime's own errors --- a wrong number of
arguments, no applicable method, a type error, integer overflow,
division by zero, an unbound variable --- are signalled as Dylan
conditions too, \code{<simple-error>} or \code{<type-error>}, and can be
handled.

There is no debugger to fall into when a condition is not handled.
Instead, the default handler (\code{default-handler}) acts as follows:

\begin{itemize}
\item An unhandled \textbf{serious condition} --- any \code{<error>},
  among others --- writes \code{Dylan error:} and the condition's message
  to the standard error stream, and the program exits with status~1.
  The standard output is flushed first, so everything printed before the
  error appears.
\item An unhandled \textbf{warning} writes \code{Dylan warning:} and its
  message to the standard error stream, and \code{signal} returns
  \code{\#f}; the program continues.
\item \code{break} writes \code{Dylan break:} and its message, and
  returns \code{\#f}.
\item \code{restart-query} and \code{return-query} do nothing.
\end{itemize}

\begin{dylan}
(define f (method (x) (+ x 1)))
(print (f 1))
(print (f "one"))
(print 'never)
\end{dylan}

\begin{shell}
$ bin/dylanc err1.dyl && ./err1; echo "status $?"
2
Dylan error: no applicable method for binary+
status 1
\end{shell}

The message names the generic function \code{binary+}: a call of
\code{+} with two arguments calls the library's binary addition directly
(Chapter~\ref{ch:internals}). In a thread other than the main one, an
unhandled serious condition ends only that thread; see
Chapter~\ref{ch:extensions}.

\begin{warnbox}
A recursion that is \emph{not} in tail position uses C stack, and a
program that recurses too deeply --- millions of frames --- overflows the
stack and is killed by the operating system with a segmentation fault,
not a Dylan error. Output it had printed but not yet flushed is lost.
Calls in tail position never do this: they run in constant space however
deep they go. Write deep recursions as tail calls, or raise the stack
limit with \cmd{ulimit -s}.
\end{warnbox}

\section{What a compiled program cannot do}\label{sec:dylanc-static}

\index{static program}
The compiler compiles a \term{static program}: everything the program
defines is known when it is compiled, and its classes and generic
functions do not change while it runs. That is what lets the compiler
lay out every class, compile every generic function's dispatch in
advance, and call methods directly. The 1992 manual describes a dynamic
language in which a running program may add methods and slots and make
new classes, and the interactive implementation allows all of that. A
compiled program does not. The restrictions, and how each shows itself,
are these.

\paragraph{Changing classes and generic functions.}
\code{add-method}, \code{remove-method}, \code{add-slot},
\code{remove-slot} and \code{freeze-methods} cannot be used. Any use of
them --- even one that would never run --- is a compile error:

\begin{shell}
$ cat s.dyl
(define-generic-function g (x))
(add-method g (method ((x <integer>)) x))
$ bin/dylanc s.dyl
s.dyl:2:1: error: add-method cannot be used: a program's classes and
generic functions are fixed when it is compiled
\end{shell}

(The message is on one line; it is broken here to fit the page.)
Define methods with \code{define-method} instead, which adds them when
the program is compiled. A class cannot be defined twice
(``\code{<a> is defined twice (classes cannot be redefined)}'').

\paragraph{Specializers and superclasses known when compiling.}
Each specializer of a method made by \code{define-method}, and each
superclass of a class, must be something the compiler can find: a class
named directly, a variable defined as a constant class, \code{(singleton
\meta{x})}, or \code{(make <singleton> object: \meta{x})} of a constant.
A variable whose value is computed when the program runs is a compile
error:

\begin{dylan}
(define c (if (> 2 1) <integer> <string>))
(define-method m ((x c)) x)
; error: a specializer must be a class or a singleton known at compile time
\end{dylan}

\noindent while this is accepted, because \code{c} is defined as a class:

\begin{dylan}
(define c <integer>)
(define-method m ((x c)) 'yes)
(print (m 3))        ; prints yes
\end{dylan}

The owner of a slot added with \code{define-slot} may be any singleton,
but a class owner must be known when compiling.

\paragraph{Making classes and generic functions when running.}
\code{(make <class> ...)} is an error when it runs (``cannot make an
instance of this class''). \code{(make <generic-function> ...)}
answers a generic function with no methods; since no method can be
added, every call of it signals ``no applicable method''. \code{seal}
works as the manual says, and \code{make-read-only} answers its
argument.

\paragraph{Numbers.}
Integers are 63 bits: from $-2^{62}$ to $2^{62}-1$, that is from
\code{-4611686018427387904} to \code{4611686018427387903}. There are no
bignums. An integer literal outside that range is a compile error; an
operation whose result falls outside it signals an error when it runs:

\begin{dylan}
(print (* 4611686018427387903 2))
; Dylan error: integer overflow (integers are 63 bits)
\end{dylan}

The same overflow rule applies to the numerator and denominator of a
ratio. \code{<complex>} exists, as the abstract superclass of
\code{<real>}, but has no instances, and a complex literal such as
\code{\#c(1 2)} is a compile error. The interactive implementation has
both bignums and complex numbers (Chapter~\ref{ch:numbers}).

\paragraph{No evaluation at run time.}
A compiled program cannot evaluate or load code while it runs: there is
no \code{eval} and no \code{load}. Any code it runs must be in its
source when it is compiled. (Macros, which run at compile time, can
compute code; see Chapter~\ref{ch:extensions}.)

\paragraph{Undefined variables.}
A reference to a variable that nothing defines is only a warning when
compiling, as shown in Section~\ref{sec:dylanc-diagnostics}. If the
reference is reached, it is an error when the program runs:

\begin{shell}
$ ./r1
Dylan error: the variable undefined-thing@dylan-user has no value
\end{shell}

Appendix~\ref{app:differences} collects these restrictions with every
other difference from the 1992 manual.

\fussy
